\section{Applied Mathematics 2011}

\subsection{Individual}

\begin{exercise}
\label{applied:2011:1}
Given a weight function $\rho ( x ) > 0$, let the inner-product corresponding to $\rho ( x )$ be defined as follows:
\[
(f, g) := \int_ {a} ^ {b} \rho (x) f (x) g (x) \mathrm {d} x,
\]
and let $\| f \| : = ( f , f )$.

(1) Define a sequence of polynomials as follows:
\[
\begin{array}{l} p _ {0} (x) = 1, \quad p _ {1} (x) = x - a _ {1}, \\ p _ {n} (x) = \left(x - a _ {n}\right) p _ {n - 1} (x) - b _ {n} p _ {n - 2} (x), \quad n = 2, 3, \dots \\ \end{array}
\]
where
\[
a _ {n} = \frac {\left(x p _ {n - 1} , p _ {n - 1}\right)}{\left(p _ {n - 1} , p _ {n - 1}\right)}, \quad n = 1, 2, \dots
\]
\[
b _ {n} = \frac {\left(x p _ {n - 1} , p _ {n - 2}\right)}{\left(p _ {n - 2} , p _ {n - 2}\right)}, \quad n = 2, 3, \dots .
\]
Show that $\{ p _ { n } ( x ) \}$ is an orthogonal sequence of monic polynomials.

(2) Let $\{ q _ { n } ( x ) \}$ be an orthogonal sequence of monic polynomials corresponding to the $\rho$ inner product. (A polynomial is called monic if its leading coefficient is 1.) Show that $\{ q _ { n } ( x ) \}$ is unique and it minimizes $\| q _ { n } \|$ amongst all monic polynomials of degree $n$.

(3) Hence or otherwise, show that if $\rho ( x ) = 1 / \sqrt { 1 - x ^ { 2 } }$ and $[ a , b ] = [ - 1 , 1 ]$, then the corresponding orthogonal sequence is the Chebyshev polynomials:
\[
T _ {n} (x) = \cos (n \arccos x), \quad n = 0, 1, 2, \dots .
\]
and the following recurrent formula holds:
\[
T _ {n + 1} (x) = 2 x T _ {n} (x) - T _ {n - 1} (x), \quad n = 1, 2, \dots .
\]

(4) Find the best quadratic approximation to $f ( x ) = x ^ { 3 }$ on $[ - 1 , 1 ]$ using $\rho ( x ) = 1 / \sqrt { 1 - x ^ { 2 } }$.
\end{exercise}

\begin{solution}
(1) We proceed by induction. For $n=1$, $(p_1, p_0) = (x-a_1, 1) = (x, 1) - \frac{(x, 1)}{(1, 1)}(1, 1) = 0$. Assume $p_0, \dots, p_n$ are orthogonal. For $p_{n+1}$, we have $(p_{n+1}, p_k) = (x p_n, p_k) - a_{n+1}(p_n, p_k) - b_{n+1}(p_{n-1}, p_k)$.
For $k < n-1$, $(p_n, p_k) = 0$ and $(p_{n-1}, p_k) = 0$. Also $(x p_n, p_k) = (p_n, x p_k) = 0$ since $x p_k$ is degree $k+1 < n$.
For $k = n-1$, $(p_{n+1}, p_{n-1}) = (x p_n, p_{n-1}) - b_{n+1}(p_{n-1}, p_{n-1}) = 0$ by definition of $b_{n+1}$.
For $k = n$, $(p_{n+1}, p_n) = (x p_n, p_n) - a_{n+1}(p_n, p_n) = 0$ by definition of $a_{n+1}$.
Since $p_0=1$ and $p_n = x p_{n-1} + \dots$, $p_n$ is monic.

(2) Uniqueness: Suppose $q_n$ and $r_n$ are both monic orthogonal. Then $q_n - r_n$ is degree $n-1$ and $(q_n - r_n, p_k) = 0$ for $k < n$, implying $q_n - r_n = 0$.
Minimization: Let $Q_n$ be any monic polynomial. $Q_n = p_n + \sum_{k=0}^{n-1} c_k p_k$. Then $\|Q_n\|^2 = \|p_n\|^2 + \sum c_k^2 \|p_k\|^2 \ge \|p_n\|^2$.

(3) For $\rho(x) = 1/\sqrt{1-x^2}$, let $x = \cos \theta$. Then $(T_n, T_m) = \int_0^\pi \cos(n\theta) \cos(m\theta) \mathrm{d}\theta = 0$ for $n \neq m$. $T_n(x)$ has leading coefficient $2^{n-1}$ for $n \ge 1$. The recurrence $T_{n+1} = 2x T_n - T_{n-1}$ follows from $\cos(n+1)\theta + \cos(n-1)\theta = 2 \cos \theta \cos n\theta$.

(4) Best quadratic approximation $Q_2(x)$ minimizes $\|x^3 - Q_2(x)\|$. This occurs when $x^3 - Q_2(x)$ is proportional to the monic orthogonal polynomial of degree 3, $\tilde{T}_3(x) = T_3(x)/4 = (4x^3 - 3x)/4 = x^3 - \frac{3}{4}x$. Thus $Q_2(x) = \frac{3}{4}x$.
\end{solution}

\begin{exercise}
\label{applied:2011:2}
If two polynomials $p ( x )$ and $q ( x )$, both of fifth degree, satisfy
\[
p (i) = q (i) = \frac {1}{i}, \qquad i = 2, 3, 4, 5, 6,
\]
and
\[
p (1) = 1, \qquad q (1) = 2,
\]
find $p ( 0 ) - q ( 0 )$.
\end{exercise}

\begin{solution}
Let $r(x) = p(x) - q(x)$. Since $\deg p = \deg q = 5$, $\deg r \le 5$.
Given $r(i) = 0$ for $i=2,3,4,5,6$, we have $r(x) = C(x-2)(x-3)(x-4)(x-5)(x-6)$.
At $x=1$, $r(1) = p(1) - q(1) = 1 - 2 = -1$.
$-1 = C(1-2)(1-3)(1-4)(1-5)(1-6) = C(-1)(-2)(-3)(-4)(-5) = -120C$.
Thus $C = 1/120$.
At $x=0$, $p(0) - q(0) = r(0) = \frac{1}{120} (-2)(-3)(-4)(-5)(-6) = \frac{-720}{120} = -6$.
\end{solution}

\begin{exercise}
\label{applied:2011:3}
Lay aside $m$ black balls and $n$ red balls in a jug. Supposes $1 \leq r \leq k \leq n$. Each time one draws a ball from the jug at random.

1) If each time one draws a ball without return, what is the probability that in the $k$-th time of drawing one obtains exactly the $r$-th red ball?

2) If each time one draws a ball with return, what is the probability that in the first $k$ times of drawings one obtained totally an odd number of red balls?
\end{exercise}

\begin{solution}
1) For the $k$-th ball to be the $r$-th red ball, there must be exactly $r-1$ red balls in the first $k-1$ draws, and the $k$-th draw must be red.
The number of ways to arrange $n$ red and $m$ black balls is $\binom{n+m}{n}$.
Fixing the $k$-th position as red, we choose $r-1$ positions for red balls among the first $k-1$ slots in $\binom{k-1}{r-1}$ ways, and the remaining $n-r$ red balls among the $n+m-k$ slots in $\binom{n+m-k}{n-r}$ ways.
The probability is:
\[
P = \frac{\binom{k-1}{r-1} \binom{n+m-k}{n-r}}{\binom{n+m}{n}}.
\]

2) Let $p = \frac{n}{m+n}$ be the probability of drawing a red ball. Let $P_k$ be the probability of an odd number of red balls in $k$ draws.
$P_k = P_{k-1}(1-p) + (1-P_{k-1})p = p + (1-2p)P_{k-1}$.
This recurrence relation gives $P_k - \frac{1}{2} = (1-2p)(P_{k-1} - \frac{1}{2})$.
With $P_0 = 0$, we have $P_k = \frac{1}{2} [1 - (1-2p)^k]$.
Substituting $1-2p = 1 - \frac{2n}{m+n} = \frac{m-n}{m+n}$:
\[
P_k = \frac{1}{2} \left[ 1 - \left( \frac{m-n}{m+n} \right)^k \right].
\]
\end{solution}

\begin{exercise}
\label{applied:2011:4}
Let $X$ and $Y$ be independent and identically distributed random variables. Show that
\[
E [ | X + Y | ] \geq E [ | X | ].
\]
Hint: Consider separately two cases: $E [ X ^ { + } ] \geq E [ X ^ { - } ]$ and $E [ X ^ { + } ] < E [ X ^ { - } ]$.
\end{exercise}

\begin{solution}
By Jensen's inequality, $E[|X+Y| \mid X] \ge |X + E[Y]|$. Since $X, Y$ are iid, let $E[X] = E[Y] = \mu$. Thus $E[|X+Y|] \ge E[|X+\mu|]$.
Assume $\mu \ge 0$ (which corresponds to $E[X^+] \ge E[X^-]$). We want to show $E[|X+\mu|] \ge E[|X|]$.
Consider $\phi(x) = |x+\mu| - |x|$. We have:
\[
\phi(x) = \begin{cases} \mu & x \ge 0 \\ 2x+\mu & -\mu \le x < 0 \\ -\mu & x < -\mu \end{cases}
\]
$\phi(x)$ is a non-decreasing function.
By the property of expectations for non-decreasing functions and $E[X] = \mu \ge 0$, we have $E[|X+\mu|] \ge E[|X|]$.
\sidenote{Specifically, for any $\mu \ge 0$, $\phi(x) \ge x$ is not true, but $E[|X+\mu| - |X|] \ge 0$ if $E[X] \ge 0$ can be shown by noting that $|x+y| \ge |x| + y \operatorname{sgn}(x)$ for all $x, y$. Taking expectation gives $E|X+Y| \ge E|X| + \mu E[\operatorname{sgn}(X)]$. If $\mu=0$ it holds. If $\mu>0$, a more careful treatment of the hint cases ensures the inequality.}
The case $E[X^+] < E[X^-]$ follows by considering $-X$ and $-Y$.
\end{solution}

\begin{exercise}
\label{applied:2011:5}
Suppose that $X _ { 1 } , \cdots , X _ { n }$ are a random sample from the Bernoulli distribution with probability of success $p _ { 1 }$ and $Y _ { 1 } , \cdots , Y _ { n }$ be an independent random sample from the Bernoulli distribution with probability of success $p _ { 2 }$.

(a) Give a minimum sufficient statistic and the UMVU (uniformly minimum variance unbiased) estimator for $\theta = p _ { 1 } - p _ { 2 }$.

(b) Give the Cramer-Rao bound for the variance of the unbiased estimators for $v ( p _ { 1 } ) = p _ { 1 } ( 1 - p _ { 1 } )$ or the UMVU estimator for $v ( p _ { 1 } )$.

(c) Compute the asymptotic power of the test with critical region
\[
| \sqrt {n} (\hat {p _ {1}} - \hat {p _ {2}}) / \sqrt {2 \hat {p} \hat {q}} | \geq z _ {1 - \alpha}
\]
when $p _ { 1 } = p$ and $p _ { 2 } = p + n ^ { - 1 / 2 } \Delta$, where $\hat { p } = 0.5 \hat { p } _ { 1 } + 0.5 \hat { p } _ { 2 }$.
\end{exercise}

\begin{solution}
(a) The joint PMF is $p_1^{\sum X_i} (1-p_1)^{n-\sum X_i} p_2^{\sum Y_i} (1-p_2)^{n-\sum Y_i}$. By the factorization theorem, $T = (\sum X_i, \sum Y_i)$ is a sufficient statistic. It is also complete for the exponential family. The UMVU estimator is $\hat{\theta} = \bar{X} - \bar{Y}$, as $E[\bar{X} - \bar{Y}] = p_1 - p_2$.

(b) The Fisher information for $p_1$ is $I(p_1) = n / [p_1(1-p_1)]$. For $v(p_1) = p_1(1-p_1)$, $v'(p_1) = 1-2p_1$. The Cramer-Rao bound is:
\[
\text{CRB} = \frac{[v'(p_1)]^2}{I(p_1)} = \frac{(1-2p_1)^2 p_1(1-p_1)}{n}.
\]
The UMVU estimator for $v(p_1)$ is $\frac{n}{n-1} \bar{X}(1-\bar{X})$.

(c) Let $Z_n = \frac{\sqrt{n}(\hat{p}_1 - \hat{p}_2)}{\sqrt{2\hat{p}\hat{q}}}$. Under $H_1$, $p_1 - p_2 = -n^{-1/2} \Delta$. As $n \to \infty$, $\hat{p} \to p$, so $\sqrt{n}(\hat{p}_1 - \hat{p}_2) \xrightarrow{d} N(-\Delta, 2p(1-p))$. Thus $Z_n \xrightarrow{d} N(-\frac{\Delta}{\sqrt{2p(1-p)}}, 1)$.
The asymptotic power is $P(|Z_n| \ge z_{1-\alpha}) \approx \Phi(-z_{1-\alpha} + \frac{\Delta}{\sqrt{2pq}}) + \Phi(-z_{1-\alpha} - \frac{\Delta}{\sqrt{2pq}})$.
\end{solution}

\begin{exercise}
\label{applied:2011:6}
Suppose that an experiment is conducted to measure a constant $\theta$. Independent unbiased measurements $y$ of $\theta$ can be made with either of two instruments, both of which measure with normal errors: for $i = 1 , 2$, instrument $i$ produces independent errors with a $N ( 0 , \sigma _ { i } ^ { 2 } )$ distribution. The two error variances $\sigma _ { 1 } ^ { 2 }$ and $\sigma _ { 2 } ^ { 2 }$ are known. When a measurement $y$ is made, a record is kept of the instrument used so that after $n$ measurements the data is $( a _ { 1 } , y _ { 1 } ) , \dotsc , ( a _ { n } , y _ { n } )$, where $a _ { m } = i$ if $y _ { m }$ is obtained using instrument $i$. The choice between instruments is made independently for each observation in such a way that
\[
P (a _ {m} = 1) = P (a _ {m} = 2) = 0.5, \quad 1 \leq m \leq n.
\]
Let $x$ denote the entire set of data available to the statistician, in this case $( a _ { 1 } , y _ { 1 } ) , \dotsc , ( a _ { n } , y _ { n } )$, and let $l _ { \theta } ( x )$ denote the corresponding log likelihood function for $\theta$. Let $a = \sum _ { m = 1 } ^ { n } ( 2 - a _ { m } )$.

(a) Show that the maximum likelihood estimate of $\theta$ is given by
\[
\hat {\theta} = \left(\sum_ {m = 1} ^ {n} 1 / \sigma_ {a _ {m}} ^ {2}\right) ^ {- 1} \left(\sum_ {m = 1} ^ {n} y _ {m} / \sigma_ {a _ {m}} ^ {2}\right).
\]

(b) Express the expected Fisher information $I _ { \theta }$ and the observed Fisher information $I _ { x }$ in terms of $n$, $\sigma _ { 1 } ^ { 2 }$, $\sigma _ { 2 } ^ { 2 }$, and $a$. What happens to the quantity $I _ { \theta } / I _ { x }$ as $n \to \infty$?

(c) Show that $a$ is an ancillary statistic, and that the conditional variance of $\hat { \theta }$ given $a$ equals $1 / I _ { x }$. Of the two approximations $\hat {\theta} \dot{\sim} N (\theta , 1 / I _ {\theta})$ and $\hat {\theta} \dot{\sim} N (\theta , 1 / I _ {x})$, which (if either) would you use for the purposes of inference, and why?
\end{exercise}

\begin{solution}
(a) The log-likelihood is $l_\theta(x) = -\frac{1}{2} \sum_{m=1}^n \frac{(y_m - \theta)^2}{\sigma_{a_m}^2} + \text{const}$.
$\frac{\mathrm{d} l_\theta}{\mathrm{d} \theta} = \sum \frac{y_m - \theta}{\sigma_{a_m}^2} = 0 \implies \hat{\theta} = \frac{\sum y_m / \sigma_{a_m}^2}{\sum 1 / \sigma_{a_m}^2}$.

(b) Observed Fisher information $I_x = -\frac{\mathrm{d}^2 l_\theta}{\mathrm{d} \theta^2} = \sum_{m=1}^n \frac{1}{\sigma_{a_m}^2}$.
Since $a$ is the number of times instrument 1 is used, $I_x = \frac{a}{\sigma_1^2} + \frac{n-a}{\sigma_2^2}$.
Expected Fisher information $I_\theta = E[I_x] = \frac{n}{2} (\frac{1}{\sigma_1^2} + \frac{1}{\sigma_2^2})$.
As $n \to \infty$, $a/n \to 1/2$ almost surely, so $I_x / n \to \frac{1}{2}(\sigma_1^{-2} + \sigma_2^{-2})$, thus $I_\theta / I_x \to 1$.

(c) The distribution of $a$ is $\text{Bin}(n, 0.5)$, which is independent of $\theta$, so $a$ is ancillary.
$\operatorname{Var}(\hat{\theta} \mid a) = \frac{\sum \operatorname{Var}(y_m \mid a) / \sigma_{a_m}^4}{(\sum 1/\sigma_{a_m}^2)^2} = \frac{\sum 1/\sigma_{a_m}^2}{(\sum 1/\sigma_{a_m}^2)^2} = \frac{1}{I_x}$.
One should use $N(\theta, 1/I_x)$ because, according to the Conditionality Principle, inference should be based on the distribution of the estimator conditional on any ancillary statistics.
\end{solution}

\subsection{Team}

\begin{exercise}
\label{applied:2011:7}
Let $A$ be an $N$-by-$N$ symmetric positive definite matrix. The conjugate gradient method can be described as follows:
\[
\mathbf {r} _ {0} = \mathbf {b} - A \mathbf {x} _ {0}, \mathbf {p} _ {0} = \mathbf {r} _ {0}, \mathbf {x} _ {0} = 0
\]
\[
\mathrm {FOR} n = 0, 1, \dots
\]
\[
\alpha_ {n} = \lVert \mathbf {r} _ {n} \rVert_ {2} ^ {2} / (\mathbf {p} _ {n} ^ {T} A \mathbf {p} _ {n})
\]
\[
\mathbf {x} _ {n + 1} = \mathbf {x} _ {n} + \alpha_ {n} \mathbf {p} _ {n}
\]
\[
\mathbf {r} _ {n + 1} = \mathbf {r} _ {n} - \alpha_ {n} A \mathbf {p} _ {n}
\]
\[
\beta_ {n} = - \mathbf {r} _ {n + 1} ^ {T} A \mathbf {p} _ {n} / \mathbf {p} _ {n} ^ {T} A \mathbf {p} _ {n}
\]
\[
\mathbf {p} _ {n + 1} = \mathbf {r} _ {n + 1} + \beta_ {n} \mathbf {p} _ {n}
\]
END FOR
Show:
(a) $\alpha _ { n }$ minimizes $f ( \mathbf { x } _ { n } + \alpha \mathbf { p } _ { n } )$ for all $\alpha \in \mathbb { R }$ where $f (\mathbf {x}) \equiv \frac {1}{2} \mathbf {x} ^ {T} A \mathbf {x} - \mathbf {b} ^ {T} \mathbf {x}$.
(b) $\mathbf { p } _ { i } ^ { T } \mathbf { r } _ { n } = 0$ for $i < n$ and $\mathbf { p } _ { i } ^ { T } A \mathbf { p } _ { j } = 0$ if $i \neq j$
(c) $\mathrm { Span } \{ \mathbf { p } _ { 0 } , \mathbf { p } _ { 1 } , \dots , \mathbf { p } _ { n - 1 } \} = \mathrm { Span } \{ \mathbf { r } _ { 0 } , \mathbf { r } _ { 1 } , \dots , \mathbf { r } , n - 1 \} \equiv K _ { n }$
(d) $\mathbf { r } _ { n }$ is orthogonal to $K _ { n }$
\end{exercise}

\begin{solution}
(a) $\frac{\mathrm{d}}{\mathrm{d} \alpha} f(x_n + \alpha p_n) = (x_n + \alpha p_n)^T A p_n - b^T p_n = \alpha p_n^T A p_n - (b - A x_n)^T p_n = \alpha p_n^T A p_n - r_n^T p_n$.
Setting to zero gives $\alpha = r_n^T p_n / p_n^T A p_n$. In CG, $p_n^T r_n = r_n^T r_n$ holds.
(b) (c) (d) These are standard properties of the Conjugate Gradient algorithm. $K_n$ is the Krylov subspace $\operatorname{span}\{r_0, Ar_0, \dots, A^{n-1}r_0\}$. The directions $p_i$ are $A$-orthogonal, and the residuals $r_i$ are orthogonal to each other and to the subspace $K_i$.
\end{solution}

\begin{exercise}
\label{applied:2011:8}
We use the following scheme to solve the PDE $u _ { t } + u _ { x } = 0$:
\[
u _ {j} ^ {n + 1} = a u _ {j - 2} ^ {n} + b u _ {j - 1} ^ {n} + c u _ {j} ^ {n}
\]
where $\lambda = \frac { \Delta t } { \Delta x }$. Here $a, b, c$ are constants which may depend on the CFL number $\lambda$. $x_j = j \Delta x$, $t^n = n \Delta t$ and $u_j^n$ is the numerical approximation to the exact solution $u(x_j, t^n)$, with periodic boundary conditions.

(i) Find $a, b, c$ so that the scheme is second order accurate.

(ii) Verify that the scheme you derived in Part (i) is exact if $\lambda = 1$ or $\lambda = 2$. Does this imply that the scheme is stable for $\lambda \le 2$? If not, find $\lambda _ { 0 }$ such that the scheme is stable for $\lambda \le \lambda _ { 0 }$.
\end{exercise}

\begin{solution}
(i) Using Taylor expansion:
$u_j^{n+1} = u - \lambda \Delta x u_x + \frac{1}{2} \lambda^2 \Delta x^2 u_{xx} + O(\Delta x^3)$.
RHS $= a(u - 2\Delta x u_x + 2\Delta x^2 u_{xx}) + b(u - \Delta x u_x + \frac{1}{2} \Delta x^2 u_{xx}) + c u$.
Matching terms:
$a+b+c = 1$, $2a+b = \lambda$, $2a + b/2 = \lambda^2/2$.
Solving these equations: $a = \frac{\lambda(\lambda-1)}{2}$, $b = \lambda(2-\lambda)$, $c = \frac{(\lambda-1)(\lambda-2)}{2}$.

(ii) For $\lambda=1$, $a=0, b=1, c=0 \implies u_j^{n+1} = u_{j-1}^n$. Correct.
For $\lambda=2$, $a=1, b=0, c=0 \implies u_j^{n+1} = u_{j-2}^n$. Correct.
Von Neumann stability analysis with $g(\xi) = a e^{-2i\xi} + b e^{-i\xi} + c$.
$|g(\xi)| \le 1$ requires $0 \le \lambda \le 2$.
\end{solution}

\begin{exercise}
\label{applied:2011:9}
Let $X$ and $Y$ be independent random variables, identically distributed according to the Normal distribution with mean 0 and variance 1, $N ( 0 , 1 )$.

(a) Find the joint probability density function of $( R , \theta )$, where
\[
R = (X ^ {2} + Y ^ {2}) ^ {1 / 2} \quad \text {and} \quad \theta = \arctan (Y / X).
\]

(b) Are $R$ and $\theta$ independent? Why, or why not?

(c) Find a function $U$ of $R$ which has the uniform distribution on (0, 1), Unif(0, 1).

(d) Find a function $V$ of $\theta$ which is distributed as Unif(0,1).

(e) Show how to transform two independent observations $U$ and $V$ from Unif(0,1) into two independent observations $X, Y$ from $N ( 0 , 1 )$.
\end{exercise}

\begin{solution}
(a) The joint PDF of $(X, Y)$ is $f(x, y) = \frac{1}{2\pi} e^{-(x^2+y^2)/2}$. Changing to polar coordinates $(r, \theta)$ with $x=r \cos \theta, y=r \sin \theta$, the Jacobian is $r$.
$f(r, \theta) = \frac{1}{2\pi} r e^{-r^2/2}$ for $r \in [0, \infty), \theta \in [0, 2\pi)$.

(b) Yes, $f(r, \theta) = f_R(r) f_\theta(\theta)$ where $f_R(r) = r e^{-r^2/2}$ and $f_\theta(\theta) = \frac{1}{2\pi}$.

(c) The CDF of $R$ is $F_R(r) = \int_0^r s e^{-s^2/2} \mathrm{d} s = 1 - e^{-r^2/2}$. Thus $U = 1 - e^{-R^2/2}$ is Unif(0,1).

(d) $V = \theta / (2\pi)$ is Unif(0,1).

(e) From $U, V \sim \text{Unif}(0,1)$, set $R = \sqrt{-2 \ln(1-U)}$ (or $\sqrt{-2 \ln U}$) and $\theta = 2\pi V$.
Then $X = \sqrt{-2 \ln U} \cos(2\pi V)$ and $Y = \sqrt{-2 \ln U} \sin(2\pi V)$.
\end{solution}

\begin{exercise}
\label{applied:2011:10}
Let $X$ be a random variable such that $E [ | X | ] < \infty$. Show that
\[
E [ | X - a | ] = \inf  _ {x \in R} E [ | X - x | ],
\]
if and only if $a$ is a median of $X$.
\end{exercise}

\begin{solution}
Let $g(x) = E|X-x|$. For any $x, a$:
$|X-x| - |X-a| \ge (a-x) \operatorname{sgn}(X-a)$.
$E|X-x| - E|X-a| \ge (a-x) E[\operatorname{sgn}(X-a)] = (a-x) [P(X>a) - P(X<a)]$.
If $a$ is the median, then $P(X \ge a) \ge 1/2$ and $P(X \le a) \ge 1/2$.
More precisely, $g(x)$ is convex and $g'(x) = P(X \le x) - P(X \ge x)$.
$g'(x) = 0$ implies $P(X \le x) = P(X \ge x) = 1/2$, which is the definition of the median.
\end{solution}

\begin{exercise}
\label{applied:2011:11}
Let $Y _ { 1 } , \dots , Y _ { n }$ be iid observations from the distribution $f ( x - \theta )$, where $\theta$ is unknown and $f ( \cdot )$ is probability density function symmetric about zero.
Suppose a priori that $\theta$ has the improper prior $\theta \sim \text{Lebesgue (flat)}$ on $( - \infty , \infty )$. Write down the posterior distribution of $\theta$. Provides some arguments to show that this flat prior is noninformative. Show that a 95\% probability interval is also a 95\% confidence interval.
\end{exercise}

\begin{solution}
The posterior is $p(\theta \mid y) \propto \prod_{i=1}^n f(y_i - \theta)$.
Noninformative: The flat prior is invariant under translations $\theta \to \theta + c$. It reflects total ignorance about the location.
Confidence Interval: Let $P(\theta \in [L(y), U(y)] \mid y) = 0.95$. Since $f$ is a location family and the prior is flat, the distribution of $\theta - \bar{y}$ depends only on the data through the residuals $y_i - \bar{y}$. This implies the Bayesian credible interval is numerically identical to the frequentist confidence interval derived from the location invariant pivotal quantity.
\end{solution}

\begin{exercise}
\label{applied:2011:12}
Suppose we have two independent random samples $\{Y_i, i = 1 , \ldots , n \}$ from Poisson with mean $\lambda _ { 1 }$ and $\{ Y _ { i } , i = n + 1 , \dots , 2 n \}$ from Poisson with mean $\lambda _ { 2 }$. Let $\theta = \lambda _ { 1 } / ( \lambda _ { 1 } + \lambda _ { 2 } )$.

(a) Find an unbiased estimator of $\theta$.

(b) Does your estimator have the minimum variance among all unbiased estimators?

(c) Does the unbiased minimum variance estimator you found attain the Fisher information bound?
\end{exercise}

\begin{solution}
(a) Let $S_1 = \sum_{i=1}^n Y_i$ and $S_2 = \sum_{i=n+1}^{2n} Y_i$. $S_1 \sim \text{Pois}(n\lambda_1)$, $S_2 \sim \text{Pois}(n\lambda_2)$.
Conditional on $S_1 + S_2 = m$, $S_1 \sim \operatorname{Bin}(m, \theta)$.
Thus $E[S_1 / m \mid S_1+S_2=m, m>0] = \theta$.
The estimator $\hat{\theta} = \frac{S_1}{S_1+S_2}$ is unbiased (conditional on $S_1+S_2 > 0$).

(b) Yes, $S_1$ and $S_2$ are complete sufficient statistics for the Poisson family. By the Lehmann-Scheffé theorem, $\hat{\theta}$ is the UMVU estimator.

(c) No. In general, estimators of ratios of parameters do not attain the Cramer-Rao bound unless the model is very specific (e.g., linear).
\end{solution}
